
% --------------------
% Final exercise (call-by-name)
% --------------------
\begin{exercise}{Call-by-name equivalence}
Consider the following definition of small-step call-by-name reduction $\to_n$:
\[
\frac{}{\,app{(lam({x}{M})},{N}) \to_n M[x := N]}
\qquad
\frac{M \to_n M'}{\,app({M},{N}) \to_n app({M'},{N})}
\]
Show that $M \eval V$ if and only if $M \steps V$.
\end{exercise}

\begin{solution}
We prove both implications.

\begin{enumerate}
    \item {($\Rightarrow$)} If $M \Downarrow V$ then $M \to_n^* V$
\\
We proceed by induction on the derivation of $M \eval V$. We need to show that if a term evaluates to a value in a single "big step," there exists a sequence of atomic transitions (reflexive-transitive closure $\to_n^*$) that achieves the same result.
\begin{enumerate}
    \item \textsc{Base Case - Axiom}: 
    Assume $lam(x.M)\Downarrow lam(x.M)$. 
    By the definition of reflexive-transitive closure, every configuration is related to itself in zero steps $lam(x.M)\to_n^0 lam(x.M)$. Therefore, $lam(x.M)\to_n^* lam(x.M)$ is trivially true.
    \item \textsc{Inductive Step - Rule}: Assume $app(M_1, M_2)\Downarrow V$. By inductive hypothesis on the premises.
    \begin{enumerate}
        \item on $M_1\Downarrow lam(x.M)$ we have $M_1\to^*_n lam(x.M)$;
        \item on $M[x:=M_2]\Downarrow V$ we have $M[x:=M_2]\to^*_n V$.
    \end{enumerate}
    From the second transition rule given (congruence), we know that if $M_1$ can take steps, then the entire application can progress: $app(M_1, M_2) \to_n^* app(lam(x.M), M_2)$. 
    \\From the first transition rule ($\beta$-reduction): $app(lam(x.M), M_2) \to_n M[x:=M_2]$.
    \\ Joining the sequences by transitivity: 
 \[app(M_1, M_2) \to_n^* app(lam(x.M), M_2)\to_n M[x:=M_2]\to_n^* V\]
 We conclude that  \[app(M_1, M_2) \to_n^*  V\]

\end{enumerate}

\item {($\Leftarrow$)} If $M \to_n^* V$ then  $M \Downarrow V$ 
We proceed by induction on the length of the reduction sequence $M \to_n^k V$.
\begin{enumerate}
    \item \textsc{Base Case ($k=0$)}: Assume $M \to^0_n V$, then $M=V$. Since $V$ is a value (lambda abstraction), by the axiom  ($\Downarrow_{(val)}$) we have $V\Downarrow V$. 

     \item \textsc{Inductive Step ($k=n+1$)}: Assume $M \to_n M'\to_n^n V$. By inductive hypothesis we know $M'\Downarrow V$. To conclude $M\Downarrow V$ we show (\textbf{\textit{Lemma}}): if $M \to_n M' $ and $M'\Downarrow V$ then $M\Downarrow V$. \\ We show the (\textbf{\textit{Lemma}}) by cases on $M \to M'$:
    \begin{enumerate}
        \item \textsc{($\beta$-reduction)}: $app(lam(x.M),N) \to_n M[x:=N]$
        \\
        Assume $ M[x:=N] \Downarrow V$. 
        \\Since we know by $(\Downarrow_{val})$ that  $lam(x.M) \Downarrow lam(x.M)$ we obtain, by Big-Step rule for application, $app(lam(x.M),N) \Downarrow V$
         \item \textsc{(Structural Rule)}: $app(M,N) \to_n app(M', N)$ with $M\to_n M'$
         \\
         Assume $app(M',N) \Downarrow V$, by inversion there must be $M$ and $P$ such that $M' \Downarrow lam(x.P)$ and $P[x:=N]\Downarrow V$. By inductive hypothesis on $M\to_n M'$ and $M'\Downarrow lam(x.P)$ we obtain $M\Downarrow lam(x.P)$. 
        \\
        Now we can use Big-Step rule for application and conclude $app(M,N)\Downarrow V$.
    \end{enumerate}
\end{enumerate}

\end{enumerate}
Equivalence is demonstrated because every Big-Step derivation tree can be linearized into a sequence of Small-Step transitions and each Small-Step transition step preserves the final "meaning" of the term defined by the Big-Step relation.

\end{solution}